\VerdierPhysicalPageAnchor{0027}%
elements of $\operatorname{Ext}_{\mathcal A}^{n}(X,Y)$, the $n$th
hyper-Ext of the complexes $X$ and $Y$ (in the sense of [1]). When $X$ and
$Y$ are bounded complexes (only finitely many components are nonzero), the
homomorphisms of degree $n$ from $X$ to $Y$ can be defined by a construction
that generalizes the Yoneda construction [3]. Composition of morphisms in
$\mathbf D(\mathcal A)$ is nothing other than Yoneda composition.

\VerdierSemanticAnchor{intro-6-plan-chapter-three-spectral-objects}
In Section 4 we study the usual spectral objects. Every complex of
$\mathcal A$ gives rise to three spectral objects with values in
$\mathbf D(\mathcal A)$, which automatically produce the essential spectral
sequences. Moreover, every sequence
$\cdots\to X_n\to X_{n+1}\to\cdots$ of complexes and morphisms of
complexes (where the transition morphisms need not be monomorphisms) gives
rise to a spectral object.

\VerdierSemanticAnchor{intro-7}
\noindent\textbf{7.}\quad
In summary, the task is to lay the foundations of a formalism, with all the
drawbacks entailed by such work: enthymemes and sorites, proofs involving no
real mathematical difficulty but often requiring chains of verifications
that are sometimes (respectively, always) tedious. This work is therefore of
interest only insofar as this formalism, through its flexibility and
generality, permits the formulation and proof of “genuine” theorems. The
concern to state and prove cohomological duality theorems was the first
motivation for this formalism. It is known that, in different contexts, there
are formally very similar theories of cohomological duality:

\VerdierSemanticAnchor{intro-7-duality-theories}
\begin{enumerate}
\renewcommand{\labelenumi}{\arabic{enumi})}
\VerdierSemanticAnchor{intro-7-duality-1}
\item The theory of duality in the cohomology of coherent sheaves on schemes
[4].

\VerdierSemanticAnchor{intro-7-duality-2}
\item The theory of duality in the cohomology of sheaves for the étale
topology of schemes [5], [6].

These theories are due to A. Grothendieck.

\VerdierSemanticAnchor{intro-7-duality-3}
\item The theory of duality in the cohomology of Galois modules, due to
J. Tate and J.-L. Verdier [7].

\VerdierSemanticAnchor{intro-7-duality-4}
\item The theory of duality in the cohomology of sheaves on a locally compact
space [9], [10].

\VerdierSemanticAnchor{intro-7-duality-5}
\item The theory of duality in the cohomology of coherent analytic sheaves.
This theory is still essentially conjectural.

\VerdierSemanticAnchor{intro-7-duality-6}
\item The theory of duality for local fields, number fields, or function
fields, and its relations with class field theory. This theory remains to be
developed. The first steps are due to M. Artin [8].
\end{enumerate}

\VerdierSemanticAnchor{intro-7-adjunction}
All these theories have as their starting point (or endpoint) an adjunction
property relating certain functors between derived categories. The

